%% =====================================================================
%% BLG 475E Software Quality and Testing - Term Project
%% Final Submission Report (Phase 1 + Phase 2)
%% IEEE Journal Template (per assignment brief)
%% Authors: Mustafa Eren KOC (150190805), Onat Baris Ercan (150210075)
%% =====================================================================
\documentclass[journal]{IEEEtran}

\usepackage{amsmath}
\usepackage{amssymb}
\usepackage{graphicx}
\usepackage{booktabs}
\usepackage{multirow}
\usepackage{array}
\usepackage{url}
\usepackage{listings}
\usepackage{xcolor}
\usepackage{hyperref}

\hypersetup{
    colorlinks=true,
    linkcolor=blue,
    citecolor=blue,
    urlcolor=cyan,
    pdftitle={BLG 475E Term Project - Final Submission Report}
}

% --- Code listing style (compact for two-column IEEE layout) ---
\definecolor{codegreen}{rgb}{0,0.6,0}
\definecolor{codegray}{rgb}{0.5,0.5,0.5}
\definecolor{codepurple}{rgb}{0.58,0,0.82}
\definecolor{backcolour}{rgb}{0.96,0.96,0.92}

\lstdefinestyle{ieeeJava}{
    backgroundcolor=\color{backcolour},
    commentstyle=\color{codegreen},
    keywordstyle=\color{magenta},
    numberstyle=\tiny\color{codegray},
    stringstyle=\color{codepurple},
    basicstyle=\ttfamily\scriptsize,
    breakatwhitespace=false,
    breaklines=true,
    captionpos=b,
    keepspaces=true,
    numbers=left,
    numbersep=5pt,
    showspaces=false,
    showstringspaces=false,
    showtabs=false,
    tabsize=2,
    columns=fullflexible,
    xleftmargin=10pt,
    aboveskip=4pt,
    belowskip=4pt
}
\lstset{style=ieeeJava}

\hyphenation{op-tical net-works semi-conduc-tor Hu-man-Eval Java-Script Power-Shell Book-Scan}

\begin{document}

\title{LLM-Based Code and Test Generation for Java:\\
From HumanEval-X Warm-Up to a Composite \texttt{BookScan} Integration Suite}

\author{Mustafa~Eren~Ko\c{c}~(150190805),~and~Onat~Bar{\i}\c{s}~Ercan~(150210075)% <-this % stops a space
\thanks{The authors are with the Department of Computer Engineering,
Istanbul Technical University, Istanbul, T\"urkiye, e-mail:
\{koc20, ercano21\}@itu.edu.tr.}% <-this % stops a space
\thanks{This is the final submission report (Phase~1 + Phase~2) of the
BLG~475E Software Quality and Testing term project (2025--2026 Spring
term), supervised by Dr.~Ay\c{s}e~Tosun~K\"uhn. The authors form a
two-person group with the explicit acknowledgement of the course
instructor. Source code, scripts, interaction logs and CI artefacts are
publicly available at
\url{https://github.com/mustafaerenkoc44/BLG-475E-Term-Project}.}}

\markboth{BLG~475E Software Quality and Testing, Final Submission Report, Spring 2026}%
{Ko\c{c} and Ercan: LLM-Based Code and Test Generation for Java}

\maketitle

\begin{abstract}
This is the final submission report for the BLG~475E Software Quality and
Testing term project. It consolidates Phase~1 (warm-up) and Phase~2
(extension) into a single coherent paper. The project investigates how
two publicly available, code-oriented Large Language Models (LLMs)
--- \texttt{Qwen2.5-Coder-1.5B-Instruct} and
\texttt{DeepSeek-Coder-1.3B-Instruct} --- behave when asked to generate
Java code and tests, first on isolated HumanEval-X prompts (Phase~1) and
then on a composite class \texttt{BookScan} that integrates three of those
prompts (Phase~2). Phase~1 establishes the methodology, the toolchain and
the baseline measurements: both LLMs reach $30/30$ base-test success on a
balanced 30-prompt subset, the improved suites raise aggregate JaCoCo
branch coverage from $96.09\,\%$ / $98.46\,\%$ (Qwen / DeepSeek) on the
base suites to $98.44\,\%$ / $100.00\,\%$, and every (task, model) pair
carries at least one hand-crafted \texttt{improvedMutation\ldots} JUnit
method. Phase~2 lifts the helpers behind Java/18 (\texttt{howManyTimes}),
Java/23 (\texttt{strlen}) and Java/27 (\texttt{flipCase}) into a single
composite \texttt{BookScan} class and compares four LLM-generated
variants (two prompt strategies $\times$ two models) against a manually
selected final implementation. The original combined prompt yields
$6/18$ (Qwen) and $8/18$ (DeepSeek) test passes; the edited combined
prompt that adds eight integration contracts lifts both models to
$17/18$; the manually selected final implementation passes $18/18$,
reaches $72/72=100\,\%$ branch and $110/110=100\,\%$ line coverage under
JaCoCo, and records a $79/93=84.95\,\%$ mutation score under PITest. The
most important Phase~2 finding is that, at the integration level, prompt
clarity does more work than model selection.
\end{abstract}

\begin{IEEEkeywords}
Large language models, integration testing, JUnit, JaCoCo, mutation
testing, PITest, prompt engineering, HumanEval-X, software quality.
\end{IEEEkeywords}

% =====================================================================
\section{Introduction}\label{sec:intro}
% =====================================================================
\IEEEPARstart{L}{arge} Language Models (LLMs) are increasingly used to
draft both source code and unit tests~\cite{chatunitest,metatestgen,
yang2024,panta2024,lit_review}. The economic appeal is obvious --- a model
can produce compilable code and candidate tests in seconds --- but the
engineering risk is equally clear: generated code can compile and pass a
happy-path test without actually meeting the specification, and generated
tests can execute without exercising any meaningful behaviour. BLG~475E
\emph{Software Quality and Testing} asks students to investigate both
sides of that trade-off on a controlled Java corpus.

This paper is the final submission report. It folds the Phase~1 warm-up
report and the Phase~2 extension report into a single paper. Phase~1
establishes the methodology, the toolchain and the baseline
measurements; Phase~2 builds on these foundations to study integration
testing of a composite class derived from three Phase~1 prompts.

The contribution of this paper is threefold. First, we show that an
orthodox semi-agentic pipeline driven by PowerShell, Python,
\texttt{llama.cpp}, JUnit~6 and JaCoCo is sufficient to reproduce every
experiment on a student laptop and on a GitHub Actions runner. Second, we
add a hand-crafted \emph{mutation-kill matrix} that complements branch
coverage in Phase~1 and pair it with PITest-driven automated mutation
analysis on the composite class in Phase~2. Third, we provide a
controlled prompt-comparison experiment that demonstrates --- with raw
JUnit logs --- that, at the integration level, explicit prompt contracts
close many more failing tests than switching between the two underlying
models does.

The remainder of the paper is organised as follows.
Section~\ref{sec:phase1-method} summarises the Phase~1 methodology, and
Section~\ref{sec:phase1-results} the Phase~1 measured results.
Section~\ref{sec:phase2-method} introduces the \texttt{BookScan} class,
the Maven build and the original-vs-edited prompt comparison.
Section~\ref{sec:phase2-challenges} documents the Phase~2 technical
challenges. Section~\ref{sec:phase2-results} reports the Phase~2
measured results: prompt-comparison outcomes, JaCoCo coverage and PITest
mutation score. Section~\ref{sec:discuss} discusses both phases together.
Section~\ref{sec:lit} contains the five-paper literature review.
Section~\ref{sec:conclusion} concludes the project.

% =====================================================================
\section{Phase 1 Methodology}\label{sec:phase1-method}
% =====================================================================

\subsection{LLM Selection and Workflow}

We adopted a \emph{semi-agentic} workflow because every transition between
steps remains observable, which makes the experiment reproducible. The
two selected models satisfy five constraints (public; code-oriented;
locally deployable in $<$8\,GB; instruction-tuned; distinct lineage):
\texttt{Qwen/Qwen2.5-Coder-1.5B-Instruct} and
\texttt{deepseek-ai/deepseek-coder-1.3b-instruct}, both deployed as GGUF
checkpoints under the same local \texttt{llama.cpp} runtime with identical
sampling parameters. The Phase~1 pipeline is driven by PowerShell~7 and
Python~3.11 scripts and re-runs end-to-end on every push via GitHub
Actions (\texttt{phase1-ci.yml}).

\subsection{Prompt Selection and Difficulty Distribution}

We selected exactly 30 prompts from the HumanEval-X Java corpus, balanced
as 10~easy / 10~moderate / 10~hard. Three of the selected prompts
(Java/18 \texttt{howManyTimes}, Java/23 \texttt{strlen}, Java/27
\texttt{flipCase}) are explicitly retained because Phase~2 lifts their
helpers into the composite \texttt{BookScan} class. The full task list is
reproduced in Appendix~\ref{app:prompts}.

\subsection{Code Generation and Base Test Conversion}

For each prompt the verbatim HumanEval-X declaration is sent to the
model. The course brief forbids editing the prompt at the generation
stage, so the only post-processing is structural normalization (Markdown
fences, prompt echo, malformed fragments). The dataset's reference test
is re-encoded as a JUnit~6 \texttt{DatasetBaseTest.java} so that JaCoCo
can be attached.

\begin{lstlisting}[language=Java, caption=Required header block prepended to every generated solution and test file.]
/* @Authors
 * Student Names: Mustafa Eren KOC, Onat Baris ERCAN
 * Student IDs: 150190805, 150210075
 */
\end{lstlisting}

\subsection{Test Improvement and Hand-Crafted Mutation Guards}

For every (task, model) pair an improved JUnit suite
\texttt{DatasetImprovedTest.java} is authored on top of the base test.
The improvement comes from three orthogonal sources: (i)~a test-smell
audit against the Phase~1 checklist that catches happy-path-only suites,
conditional logic in tests, missing boundary cases and assertion
roulette; (ii)~an equivalence-partitioning document with valid classes,
invalid classes and boundary conditions, where rows that are not already
covered by the base test become the driver for at least one new
\texttt{@Test} method; (iii)~hand-crafted mutation guards. Each improved
suite has at least one \texttt{improvedMutation\ldots} method that fails
if a specific operator-level mutation (Table~\ref{tab:mutation-catalogue})
is introduced into the solution.

\begin{table}[!t]
\renewcommand{\arraystretch}{1.15}
\caption{Hand-crafted mutation operators exercised in Phase 1.}
\label{tab:mutation-catalogue}
\centering
\footnotesize
\begin{tabular}{@{}lp{1.9cm}p{3.5cm}@{}}
\toprule
\textbf{Op} & \textbf{Family} & \textbf{Representative use} \\
\midrule
ROR & Relational op.        & Java/3 \verb|< 0| vs \verb|<= 0| \\
AOR & Arithmetic op.        & Java/2 \texttt{Math.floor} vs \texttt{(int)} \\
LCR & Logical connector     & Java/9 \verb|null  ||  size==0| \\
UOI & Unary insertion       & Java/31 negative-prime guard \\
BND & Boundary literal      & Java/25 / 36 numeric guards \\
SBR & Statement removal     & Java/0 self-pair skip \\
RV  & Return-value swap     & Java/39 non-positive index \\
NPE & Null replacement      & every I1 invalid class \\
EXC & Exception flip        & Java/11 unequal-length inputs \\
\bottomrule
\end{tabular}
\end{table}

\subsection{Refactoring Loop}

If a normalized solution fails the dataset base test, the workflow does
not silently fix the code: it issues a refactoring prompt to the same
model, attaches the failure log and the broken solution, and stores the
resulting patch as the new solution. Seven tasks (Java/10, /12, /17, /19,
/23, /27, /39) needed exactly one refactoring round before reaching
base-test success on at least one model; all seven were eventually fixed
and now compile and pass on both models.

% =====================================================================
\section{Phase 1 Results}\label{sec:phase1-results}
% =====================================================================

\subsection{Code Generation Outcomes}

After the refactoring loop, both models compile and pass their base
tests on all 30 prompts (Table~\ref{tab:codegen-final}).

\begin{table}[!t]
\renewcommand{\arraystretch}{1.15}
\caption{Phase 1 code generation outcomes after the refactoring loop.}
\label{tab:codegen-final}
\centering
\footnotesize
\begin{tabular}{@{}lcc@{}}
\toprule
\textbf{Model} & \textbf{Compiled} & \textbf{Base passed} \\
\midrule
Qwen2.5-Coder-1.5B-Instruct-GGUF    & 30 / 30 & 30 / 30 \\
DeepSeek-Coder-1.3B-Instruct-GGUF   & 30 / 30 & 30 / 30 \\
\bottomrule
\end{tabular}
\end{table}

\subsection{Branch Coverage}

\begin{table}[!t]
\renewcommand{\arraystretch}{1.15}
\caption{Phase 1 aggregate branch coverage (base versus improved).}
\label{tab:phase1-cov}
\centering
\footnotesize
\begin{tabular}{@{}lcccc@{}}
\toprule
& \multicolumn{2}{c}{\textbf{Base suite}} & \multicolumn{2}{c}{\textbf{Improved suite}} \\
\cmidrule(lr){2-3}\cmidrule(lr){4-5}
\textbf{Model}    & \textbf{Cov./Tot.} & \textbf{\%} & \textbf{Cov./Tot.} & \textbf{\%} \\
\midrule
Qwen 1.5B     & 123 / 128 & 96.09 & 126 / 128 &  98.44 \\
DeepSeek 1.3B & 128 / 130 & 98.46 & 130 / 130 & 100.00 \\
\bottomrule
\end{tabular}
\end{table}

The two remaining Qwen-side misses sit on Java/001 \texttt{separateParenGroups}
and Java/006 \texttt{parseNestedParens}, where the Qwen-generated parser
contains a private helper-exit branch that is not reachable through the
public API under the dataset contract. We document these residual misses
in Section~\ref{sec:discuss}.

\subsection{Hand-Crafted Mutation-Kill Status}

By construction, the improved suites kill $100\,\%$ of the hand-crafted
Phase~1 mutants and the LLM-produced base suites kill $0\,\%$ of them.
Phase~2 supplements this with PITest measurements on the composite
\texttt{BookScan} class.

\subsection{Model-Specific Behavioural Divergence}

For several Phase~1 tasks the two models produced semantically distinct
implementations (Table~\ref{tab:divergence}); the improved suites carry
different assertions for those tasks so the suite documents what each LLM
actually did.

\begin{table*}[!t]
\renewcommand{\arraystretch}{1.15}
\caption{Phase 1 tasks where Qwen and DeepSeek diverge on observable behaviour.}
\label{tab:divergence}
\centering
\footnotesize
\begin{tabular}{@{}lp{14.0cm}@{}}
\toprule
\textbf{Task} & \textbf{Behavioural disagreement} \\
\midrule
Java/2  & \texttt{Math.floor} (Qwen) vs \texttt{(int)} cast (DeepSeek) for negative input. \\
Java/4  & \texttt{Optional.orElse(0.0)} (Qwen) vs division-by-zero \(\to\) \texttt{NaN} (DeepSeek) on empty list. \\
Java/19 & \texttt{indexOf == -1} (Qwen) vs \texttt{HashMap.get == null} \(\to\) NPE (DeepSeek). \\
Java/20 & First tied pair (Qwen) vs accumulate-all-ties (DeepSeek). \\
Java/21 & Returns a fresh list (Qwen) vs mutates input in place (DeepSeek). \\
\bottomrule
\end{tabular}
\end{table*}

% =====================================================================
\section{Phase 2 Methodology}\label{sec:phase2-method}
% =====================================================================

\subsection{The \texttt{BookScan} Class}

Phase~2 lifts three Phase~1 helpers (Java/18 \texttt{howManyTimes}, Java/23
\texttt{strlen}, Java/27 \texttt{flipCase}) into a single composite class
\texttt{BookScan} that must (i)~count how many words of a requested length
appear in a multi-line text and on which lines they appear; (ii)~count how
many times a requested whole word appears across lines, case-insensitively;
and (iii)~preserve the Phase~1 helper semantics so that the three helpers
remain individually callable and individually testable.

\begin{lstlisting}[language=Java, caption=Public surface of the final BookScan class.]
public record ScanResult(
    int targetLength,
    int totalOccurrences,
    List<Integer> matchingLines,
    Map<Integer,Integer> occurrencesByLine,
    Map<Integer,List<String>> matchedWordsByLine) { }

public record WordSearchResult(
    String word,
    int totalOccurrences,
    List<Integer> matchingLines,
    Map<Integer,Integer> occurrencesByLine) { }

public ScanResult scanByWordLength(String text, int targetLength);
public WordSearchResult scanWord(String text, String word);
public int    howManyTimes(String string, String substring);
public int    strlen      (String string);
public String flipCase    (String string);
\end{lstlisting}

The tokenisation regex is \verb|[A-Za-z0-9]+(?:['-][A-Za-z0-9]+)*|, which
keeps apostrophes and hyphens \emph{inside} a token (so that
\texttt{can't} and \texttt{co-op} count as one word each) but rejects
them as leading or trailing characters. Whole-word matching is
implemented by canonicalising each line into a space-separated list of
lowercase tokens and then padding the needle with spaces, so that
partial matches inside longer words (e.g.\ ``echo'' inside ``echoic'')
cannot fire.

Listing~\ref{lst:scanlength} reproduces the core
\texttt{scanByWordLength} method, which illustrates the three-helper
integration: \texttt{tokenizeWords} delivers the regex-driven token
list; the loop accumulates per-line counters and matched words; the
result is wrapped in an immutable \texttt{ScanResult} record.

\begin{lstlisting}[language=Java, label=lst:scanlength,
   caption=Core scanByWordLength method (essential shape).]
public ScanResult scanByWordLength(
        String text, int targetLength) {
    if (text == null || text.isBlank()
            || targetLength <= 0) {
        return emptyScanResult(targetLength);
    }
    List<Integer> matchingLines  = new ArrayList<>();
    Map<Integer,Integer> byLine  = new LinkedHashMap<>();
    Map<Integer,List<String>> ws = new LinkedHashMap<>();
    int total = 0;
    String[] lines = text.split("\\R", -1);
    for (int i = 0; i < lines.length; i++) {
        int ln = i + 1;
        List<String> hits = new ArrayList<>();
        for (String tok : tokenizeWords(lines[i])) {
            if (strlen(tok) == targetLength) hits.add(tok);
        }
        if (!hits.isEmpty()) {
            matchingLines.add(ln);
            byLine.put(ln, hits.size());
            ws.put(ln, List.copyOf(hits));
            total += hits.size();
        }
    }
    return new ScanResult(targetLength, total,
        List.copyOf(matchingLines),
        Map.copyOf(byLine), Map.copyOf(ws));
}
\end{lstlisting}

\subsection{Maven + Surefire + JaCoCo + PITest Pipeline}

Phase~2 replaces the Phase~1 PowerShell+Python harness with a Maven build
for the production code: JDK~21 (Temurin) with Java release target~21,
JUnit~6 (\texttt{junit-jupiter} 6.0.3) on \texttt{maven-surefire-plugin}
3.5.3, JaCoCo~0.8.12 via \texttt{jacoco-maven-plugin} (prepare-agent +
report), and PITest~1.20.5 with \texttt{pitest-junit5-plugin}~1.2.3
exposed through an opt-in Maven profile \texttt{mutation} so that the
default \texttt{mvn clean verify} stays fast. The PowerShell + Python
driver is retained as the prompt-comparison harness
(\texttt{run\_phase2\_prompt\_comparison.py}). GitHub Actions re-runs the
Phase~2 build on every push (\texttt{phase2-ci.yml}).

\subsection{Original-vs-Edited Prompt Strategies}

The headline experimental question of Phase~2 is: ``once the prompt has
to specify a whole class rather than a single function, which matters
more --- the underlying model, or the structure of the prompt itself?'' We
answer it by running the same two models against two prompt strategies,
plus a manually selected final implementation, all evaluated against the
same JUnit suite.

\subsubsection{Original Combined Prompt}
The original combined prompt simply concatenates the three HumanEval-X
task descriptions and asks the model to build a class around them.

\subsubsection{Edited Combined Prompt}
The edited combined prompt adds the integration contracts that the
original prompt leaves implicit (Table~\ref{tab:prompt-edits}); these are
the contracts the integration suite would otherwise have to discover by
trial and error.

\begin{table}[!t]
\renewcommand{\arraystretch}{1.15}
\caption{Integration contracts added in the edited combined prompt.}
\label{tab:prompt-edits}
\centering
\footnotesize
\begin{tabular}{@{}p{2.6cm}p{4.6cm}@{}}
\toprule
\textbf{Topic} & \textbf{Added requirement} \\
\midrule
Line numbering    & one-based, each matching line listed at most once \\
Tokenisation      & apostrophes and hyphens stay inside a single token \\
Whole-word search & must not match inside a longer word \\
Case policy       & search is case-insensitive; uppercase-only queries canonicalised \\
Query hygiene     & trim the query, reject blank or multi-token queries \\
Invalid inputs    & null or blank text returns empty-result records \\
Result shape      & immutable \texttt{ScanResult} / \texttt{WordSearchResult} \\
Helper reuse      & \texttt{howManyTimes}, \texttt{strlen}, \texttt{flipCase} remain callable \\
\bottomrule
\end{tabular}
\end{table}

\subsubsection{Manually Selected Final Implementation}
Neither edited-prompt variant clears the full strengthened suite (each
loses exactly one helper-hardening regression test). The repository
therefore also stores a manually selected final implementation
(\texttt{Phase2/src/main/java/BookScan.java}) that closes the last gap.
The Maven build, the JaCoCo report and the PITest profile target this
implementation.

\subsection{Integration and Regression Test Suites}

The Phase~2 suite contains 12 integration tests
(\texttt{BookScanIntegrationTest}) and 6 focused regression / hardening
tests (\texttt{BookScanRegressionTest}), 18 tests in total, shared by all
five configurations. Tables~\ref{tab:int-tests} and \ref{tab:reg-tests}
enumerate each test method together with the integration contract or
helper semantic it exercises.

\begin{table*}[!t]
\renewcommand{\arraystretch}{1.15}
\caption{Integration tests in BookScanIntegrationTest.}
\label{tab:int-tests}
\centering
\footnotesize
\begin{tabular}{@{}p{6.0cm}p{10.0cm}@{}}
\toprule
\textbf{Test method} & \textbf{Integration contract} \\
\midrule
\texttt{scanByWordLengthAggregates}\-\texttt{CountsAndLinesAcrossText} & multi-line aggregation, per-line matched words \\
\texttt{scanByWordLengthTreats}\-\texttt{ApostrophesAndHyphensAs}\-\texttt{SingleWords} & token-boundary preservation \\
\texttt{scanByWordLengthReturns}\-\texttt{EmptyResultForInvalidInputs} & null / blank / non-positive length \\
\texttt{scanByWordLengthKeeps}\-\texttt{MatchingLinesUniqueEvenWith}\-\texttt{ManyHits} & line-list uniqueness under repeated hits \\
\texttt{scanWordCountsWholeWord}\-\texttt{MatchesIgnoringCaseAcrossLines} & case-insensitive whole-word matching \\
\texttt{scanWordFindsTrimmed}\-\texttt{QueriesNextToPunctuation} & trimmed queries, punctuation handling \\
\texttt{scanWordRejectsBlankOr}\-\texttt{MultiTokenQueries} & query validation \\
\texttt{scanWordReturnsEmpty}\-\texttt{ResultForNullInputsAnd}\-\texttt{BlankTexts} & null / blank inputs produce stable empty results \\
\texttt{scanWordDoesNotMatch}\-\texttt{InsideLongerWords} & partial-match rejection (e.g.\ ``echo'' inside ``echoic'') \\
\texttt{scanWordCanonicalizes}\-\texttt{UppercaseQueriesBefore}\-\texttt{WholeWordMatching} & uppercase queries normalised via \texttt{flipCase} \\
\texttt{scanWordHandlesAlphanumeric}\-\texttt{WholeWordQueries} & alphanumeric whole-word tokens \\
\texttt{scanWordNormalizesMixed}\-\texttt{CaseTokensWithoutFlipArtifacts} & mixed-case normalisation remains stable \\
\bottomrule
\end{tabular}
\end{table*}

\begin{table*}[!t]
\renewcommand{\arraystretch}{1.15}
\caption{Regression / hardening tests in BookScanRegressionTest.}
\label{tab:reg-tests}
\centering
\footnotesize
\begin{tabular}{@{}p{6.0cm}p{10.0cm}@{}}
\toprule
\textbf{Test method} & \textbf{Phase 1 helper under test} \\
\midrule
\texttt{howManyTimesCounts}\-\texttt{OverlappingSubstrings} & Java/18 overlap semantics \\
\texttt{howManyTimesReturnsZeroFor}\-\texttt{EmptyNullOrOversizedNeedles} & Java/18 guards on blank / oversized needles \\
\texttt{flipCasePreservesSymbols}\-\texttt{WhileTogglingLetters} & Java/27 case flip on symbols and digits \\
\texttt{strlenThrowsOnNullAnd}\-\texttt{ReturnsLengthOtherwise} & Java/23 null contract and length \\
\texttt{tokenizeWordsReturnsFresh}\-\texttt{MutableEmptyListForNullInput} & helper hardening, null tokenisation \\
\texttt{privateCanonicalizeWord}\-\texttt{RejectsBlankTokens} & helper hardening, blank-token collapse \\
\bottomrule
\end{tabular}
\end{table*}

The integration tests are assembled from a cross-product of three sources:
the Phase~1 test-smell checklist, the Phase~2 black-box assessment under
\texttt{Phase2/docs/analysis/black\_box\_assessment.md}, and the concrete
failure signatures we observed while running the original-prompt variants.
Every test asserts at least one structural property (counts) and at least
one shape property (matching-line list or per-line map), so a pass
genuinely implies integration correctness, not just one scalar equality.

\subsection{Black-Box Assessment}

Each public method of \texttt{BookScan} has its own
equivalence-partitioning document under
\texttt{Phase2/docs/analysis/black\_box\_assessment.md}: valid classes,
invalid classes, and boundary conditions. Every public class is mapped to
at least one executed JUnit method. Table~\ref{tab:bb-scanwordlength}
reproduces the equivalence classes for the most behaviour-rich method,
\texttt{scanByWordLength}, as a representative example.

\begin{table}[!t]
\renewcommand{\arraystretch}{1.15}
\caption{Black-box equivalence classes for \texttt{scanByWordLength}.
V$=$valid, I$=$invalid, B$=$boundary.}
\label{tab:bb-scanwordlength}
\centering
\footnotesize
\begin{tabular}{@{}lp{3.0cm}p{2.4cm}c@{}}
\toprule
\textbf{ID} & \textbf{Description} & \textbf{Representative input} & \textbf{Cov.} \\
\midrule
V1 & Multi-line text, exact-length hits & ``Echo code\textbackslash nThis here'', $4$ & yes \\
V2 & Repeated hits on same line        & ``same size size'', $4$           & yes \\
V3 & Apostrophe / hyphen tokens         & ``can't / co-op'', $5$            & yes \\
I1 & Null text                          & \texttt{null}, $4$                & yes \\
I2 & Blank text                         & ``\textbackslash n\textbackslash t'', $4$ & yes \\
I3 & Non-positive target length         & ``hello'', $0$ or $-1$            & yes \\
B1 & Empty matched-line list            & ``aa bbb cccc'', $5$              & yes \\
B2 & Single-token line                  & ``hello'', $5$                    & yes \\
B3 & Target equals max token length     & ``aa bbbbb'', $5$                 & yes \\
\bottomrule
\end{tabular}
\end{table}

The full assessment for \texttt{scanWord}, \texttt{howManyTimes},
\texttt{strlen} and \texttt{flipCase} follows the same pattern; every
row is annotated with the exact JUnit method that exercises it, so the
mapping from contract to test is auditable.

\subsection{Prompt-Comparison Harness}

The Python runner
\texttt{scripts/python/run\_phase2\_prompt\_comparison.py} is invoked
through the PowerShell wrapper
\texttt{scripts/Run-Phase2PromptComparison.ps1}. For every candidate
\texttt{BookScan.java} (four prompt variants plus the manually selected
final implementation) it copies the candidate and the shared JUnit suite
into an isolated working directory, compiles with \texttt{javac} against
the JUnit~6 console jar, runs JUnit under the JaCoCo agent, renders a
JaCoCo CSV from the resulting \texttt{jacoco.exec}, and extracts branch
coverage for the \texttt{BookScan} class. All raw artefacts
(\texttt{compile.stdout}, \texttt{compile.stderr}, \texttt{junit.stdout},
\texttt{jacoco.exec}, \texttt{jacoco.csv} and the copied source files)
are stored under
\texttt{Phase2/results/artifacts/prompt-comparison/}\,
\texttt{<strategy>/<model>}, so every figure in this report can be
reproduced from recorded data rather than from in-memory state.

% =====================================================================
\section{Phase 2 Technical Challenges}\label{sec:phase2-challenges}
% =====================================================================

\subsection{Unicode Workspace Path}
The Maven \texttt{jacoco-maven-plugin} originally placed
\texttt{jacoco.exec} under a path containing the Turkish character
\.{I} (\verb|D:\Indirilenler\...|) which produced corrupted byte streams.
The \texttt{destFile} and \texttt{dataFile} are now redirected to the
system tmpdir, the prompt-comparison harness uses a relative
\texttt{jacoco.exec}, and the Maven local repository is pinned to a
workspace-local \texttt{.m2}.

\subsection{JaCoCo / PITest Coexistence}
PITest also instruments the bytecode; running both PITest and JaCoCo's
\texttt{prepare-agent} in the same Maven invocation produced a classpath
clash. The \texttt{mutation} profile now sets \texttt{<jacoco.skip>true</jacoco.skip>}
so JaCoCo's agent does not fire when the user asks for a mutation run via
\texttt{mvn -P mutation test}. The default \texttt{mvn clean verify} keeps
JaCoCo enabled.

\subsection{Original Prompt Failure Modes}

Table~\ref{tab:failure-modes} lists the four distinct defect classes per
model that the original combined prompt produces on the JUnit suite.
Every defect class is an integration-level contract that is missing from
the original prompt and that would have been invisible at the
single-function Phase~1 level.

\begin{table*}[!t]
\renewcommand{\arraystretch}{1.15}
\caption{Original combined prompt: defect classes observed on the JUnit suite.}
\label{tab:failure-modes}
\centering
\footnotesize
\begin{tabular}{@{}lp{14.5cm}@{}}
\toprule
\textbf{Model} & \textbf{Defect classes hit by the JUnit suite} \\
\midrule
Qwen     & zero-based / non-unique line numbers; raw space-splitting destroys \texttt{can't}/\texttt{co-op}; case-sensitive substring search; missing canonicalisation for uppercase or alphanumeric queries; missing helper extraction. \\
DeepSeek & punctuation stripping changes token identity; substring matching inside longer words; over-counting in case-insensitive search; weak handling of punctuation around trimmed queries; alphanumeric over-counting; missing helper extraction. \\
\bottomrule
\end{tabular}
\end{table*}

% =====================================================================
\section{Phase 2 Results}\label{sec:phase2-results}
% =====================================================================

\subsection{Prompt-Comparison Outcomes}

Table~\ref{tab:phase2-cmp} aggregates the raw CSV in
\texttt{Phase2/results/prompt\_comparison\_results.csv}. The original
prompt lets both models compile but produces significant JUnit failures;
the edited prompt removes all public-integration failures but still leaves
one helper-hardening regression in each model; the manually selected
final implementation closes that gap.

\begin{table}[!t]
\renewcommand{\arraystretch}{1.15}
\caption{Prompt-comparison results (BookScan class only).}
\label{tab:phase2-cmp}
\centering
\footnotesize
\begin{tabular}{@{}llcc@{}}
\toprule
\textbf{Strategy} & \textbf{Model} & \textbf{Tests} & \textbf{Branch cov.} \\
\midrule
original-combined & Qwen     & 6 / 18  & n/a \\
original-combined & DeepSeek & 8 / 18  & n/a \\
edited-combined   & Qwen     & 17 / 18 & n/a \\
edited-combined   & DeepSeek & 17 / 18 & n/a \\
selected-final    & manual   & 18 / 18 & 72/72 = 100\,\% \\
\bottomrule
\end{tabular}
\end{table}

The headline finding is that the edited prompt removes all
public-integration failures for both models, while switching between
Qwen and DeepSeek inside a fixed prompt strategy moves the pass rate by
no more than two test cases. \emph{Once the integration contracts are
explicit, prompt clarity does more work than model selection.}

\subsection{Final BookScan Coverage}

The Maven JaCoCo CSV at \texttt{Phase2/target/site/jacoco/jacoco.csv}
reports for the manually selected final implementation: branch coverage
$72/72=100.00\,\%$, line coverage $110/110=100.00\,\%$, and
cyclomatic-complexity-covered $51/51=100.00\,\%$. No JaCoCo branch or
line misses remain in the final Maven report.

\subsection{PITest Mutation Score}

Running the opt-in Maven profile \texttt{mutation}
(\texttt{mvn -P mutation test}) under the ASCII-safe configuration
described in Section~\ref{sec:phase2-challenges} yields:
total mutants generated by PITest: $93$;
killed: $79$; survived: $14$; mutation score: $79/93=84.95\,\%$;
mutated-line coverage: $110/110=100.00\,\%$. The 14 surviving mutants are
concentrated on equivalent or defensive helper branches; the per-method
breakdown is given in Table~\ref{tab:pit-survivors} and the full per-line
inventory in \texttt{Phase2/docs/analysis/mutation\_assessment.md}.

\begin{table}[!t]
\renewcommand{\arraystretch}{1.15}
\caption{PITest survivors by location and mutation family.}
\label{tab:pit-survivors}
\centering
\footnotesize
\begin{tabular}{@{}lcp{4.0cm}@{}}
\toprule
\textbf{Method} & \textbf{Surv.} & \textbf{Reason class} \\
\midrule
\texttt{canonicalizeWord}    & 6 & equivalent normalisation paths \\
\texttt{scanByWordLength}    & 3 & defensive blank-input guard \\
\texttt{howManyTimes}        & 2 & defensive empty / oversized needle \\
\texttt{scanWord}            & 1 & defensive null query handling \\
\texttt{tokenizeWords}       & 1 & equivalent regex-matcher boundary \\
\texttt{flipCase}            & 1 & equivalent on non-letter character \\
\midrule
\textbf{total}               & \textbf{14} & 79 / 93 = 84.95\,\% killed \\
\bottomrule
\end{tabular}
\end{table}

Every surviving mutant has been classified as either \emph{equivalent}
(no observable behavioural change for any input the public API accepts)
or \emph{defensive} (covers a private guard whose only purpose is to
keep the public method safe against unreachable inputs). No surviving
mutant is in the public hot path of \texttt{scanByWordLength} or
\texttt{scanWord}; every public-behaviour mutant is killed.

\subsection{Helper-Hardening Regression Closure}

Two of the six regression tests are specifically designed to close the
last private-helper coverage gap that both edited-prompt variants leave
open: \texttt{tokenizeWordsReturnsFreshMutableEmptyListForNullInput}
hardens the null-input behaviour of the tokenisation helper, and
\texttt{privateCanonicalizeWordRejectsBlankTokens} pins down the
blank-token normalisation rule that distinguishes the manually selected
final implementation from the two edited-prompt variants.

% =====================================================================
\section{Discussion}\label{sec:discuss}
% =====================================================================

\subsection{Phase 1 versus Phase 2 Lessons}
The two phases tell a consistent story. In Phase~1, where each prompt is
a single function, the dataset's reference test plus the
equivalence-class extension already cover most of the meaningful
behaviour, and the choice of model determines mostly defensive coding
style. In Phase~2, where the prompt has to specify an integrated class,
the dataset reference test no longer exists, the integration contracts
are not stated by the prompt by default, and the choice of model is
dwarfed by the choice of prompt: both models reach the same $\sim\!\!94\,\%$
pass rate under the edited prompt, even though under the original prompt
they sat at $\sim\!\!33$--$\!44\,\%$.

\subsection{Prompt Engineering Outweighs Model Selection at the Integration Level}
Going from the original prompt to the edited prompt closes 11/18 (Qwen)
and 9/18 (DeepSeek) tests; switching from one model to the other inside a
fixed prompt closes only 0--2 tests. This is consistent with the recent
literature on LLM unit-test generation~\cite{metatestgen,yang2024,
panta2024,lit_review}, where explicit specifications repeatedly
outperform model upsizing.

\subsection{Hand-Crafted versus Automated Mutation Testing}
Phase~1 forbids automated mutation tooling, so the hand-crafted
\texttt{improvedMutation\ldots} guards are an executable specification
for the operator-level mutations the suite is supposed to catch.
Phase~2 uses PITest in addition, which generates mutants automatically
and removes the human-curation bias from the mutation set. The two
perspectives are complementary: PITest discovers mutants we did not
think of (and produces a measurable $84.95\,\%$ kill rate that we now
have to defend), while the hand-crafted guards force us to articulate
\emph{why} a particular mutation should be killed.

\subsection{Defect-Class to Failing-Test Mapping}

The data behind Section~\ref{sec:phase2-results} is more granular than
the headline numbers suggest. The original-prompt Qwen variant, for
example, hits 12 distinct failing tests, but those failures are caused
by only four underlying integration contracts that are missing from the
prompt: zero-based versus one-based line numbers (3 failing tests),
raw-space tokenisation that destroys \texttt{can't}/\texttt{co-op}
(3 tests), case-sensitive substring search (2 tests), and missing
canonicalisation for uppercase or alphanumeric queries (4 tests). The
edited combined prompt fixes all four contracts at once, which is why
its pass rate jumps from $6/18$ to $17/18$ rather than improving in a
linear ``one test per added clarification'' fashion. The remaining
$1/18$ gap is the helper-hardening regression
(\texttt{privateCanonicalizeWordRejectsBlankTokens}) that neither
edited-prompt variant produces a private \texttt{canonicalizeWord}
helper to satisfy. The same shape holds for DeepSeek with slightly
different per-class counts. We read this as evidence that the
\emph{edited prompt is not just better in headline pass rate, it
collapses several distinct defect classes into a single specification
that the LLM can act on}.

\subsection{Continuous Integration}

The Phase~2 GitHub Actions workflow (\texttt{phase2-ci.yml}) re-runs the
entire build on every push: it sets up JDK~21, downloads JUnit~6 and
JaCoCo~0.8.12, runs \texttt{mvn clean verify} (Surefire 18-test suite +
JaCoCo report), executes the prompt-comparison harness, summarises the
prompt-comparison CSV, asserts that the \texttt{selected-final} row stays
green, runs the \texttt{mutation} profile as a best-effort step, and
uploads everything under \texttt{Phase2/results}, \texttt{target/site/jacoco}
and \texttt{target/pit-reports} as build artefacts. The \texttt{mutation}
profile is opt-in (continue-on-error) so that a transient PITest
mismatch does not gate the rest of the pipeline; the mutation score
reported in this paper was reproduced under the ASCII-safe local
configuration.

\subsection{Threats to Validity}
\emph{Internal validity} is limited by local-model non-determinism;
mitigated by pinning the sampling seed, the model checkpoints, and by
re-running the entire pipeline under GitHub Actions on every push.
\emph{External validity} is limited because Phase~1 evaluates a
30-prompt subset of single-function puzzles and Phase~2 evaluates a
single composite class; the headline numbers should be read as a
controlled study, not as production benchmarks. \emph{Construct
validity} depends on branch coverage and mutation kill being faithful
proxies for test quality; we use both, and report PITest as the
automated counterpart to the hand-crafted Phase~1 guards. We also offset
the human-curation bias of the Phase~1 mutation set with an
independently-generated PITest mutation set in Phase~2, on a class
(\texttt{BookScan}) that was not part of the original Phase~1 corpus.

\subsection{Future Work}
Two extensions would tighten the project further. First, sweeping the
sampling seed across both models and reporting confidence intervals on
the prompt-comparison delta would distinguish ``Qwen and DeepSeek are
statistically different on the original prompt'' from ``the difference
is within sampling noise''. Second, replacing the manually selected
final implementation with an LLM-driven repair loop targeted at the
surviving PITest mutants would close the loop that Phase~2 deliberately
leaves open: today the final implementation is hand-curated, and a
fully agentic version of the same workflow could re-issue the failing
test as a refactoring prompt and converge on the same hardened helpers
without human intervention.

% =====================================================================
\section{Literature Review}\label{sec:lit}
% =====================================================================

The literature set was selected to cover five complementary perspectives
on LLM-driven test generation, all published within the last three years
(2023--2026); the selection was performed by hand using Google Scholar
and IEEE Xplore.

\textbf{ChatUniTest~\cite{chatunitest}} is an early reference for LLM-based
unit-test generation. Its main contribution is wrapping the model inside
a generation--validation--repair loop. This motivates our normalisation
plus repair loop in Phase~1 and the prompt-comparison harness in
Phase~2.

\textbf{Meta TestGen-LLM~\cite{metatestgen}} accepts model output only
after it passes build, execution, flakiness and coverage-gain filters.
Our pipeline uses the same instinct.

\textbf{Yang \emph{et al.}~\cite{yang2024}} report that prompt structure
materially affects compilation and coverage outcomes, and that
open-source models can be competitive but still need additional
engineering. Our Phase~2 prompt-comparison data mirrors this exactly:
under the original prompt both models sit between 33\,\% and 45\,\%
pass rate; under the edited prompt both reach 94\,\%.

\textbf{Panta~\cite{panta2024}} combines LLM generation with static
control-flow analysis and dynamic coverage feedback. Our improved-test
step in Phase~1 and our PITest profile in Phase~2 are simpler instances
of the same philosophy.

\textbf{Tasarsu, Tokmak and Catal~\cite{lit_review}} synthesise 38
peer-reviewed studies from 2020--2025. Their findings confirm patterns
visible in our project: coverage remains the dominant evaluation metric;
post-processing and integration logic are common; maintainability and
robustness are still weaker than simple coverage reporting. We
complement coverage with hand-crafted mutation guards in Phase~1 and
with automated PITest in Phase~2 specifically because of this
observation.

% =====================================================================
\section{Conclusion}\label{sec:conclusion}
% =====================================================================

The combined Phase~1 + Phase~2 submission shows that two small open-weight
code-LLMs can drive a reproducible Java evaluation pipeline when they are
surrounded by normalization, repair, coverage and mutation-based
guardrails. On 30 HumanEval-X Java prompts both models reach $30/30$
base-test success after at most one logged refactoring round; the
improved suites lift aggregate branch coverage from $96.09\,\%$ /
$98.46\,\%$ to $98.44\,\%$ / $100.00\,\%$; every (task, model) pair
carries at least one hand-crafted mutation guard. Phase~2 lifts the
helpers behind Java/18, Java/23 and Java/27 into a composite
\texttt{BookScan} class and runs the same two models against an original
combined prompt, an edited combined prompt and a manually selected final
implementation under a Maven build with Surefire, JaCoCo and PITest.
The original prompt produces $6/18$ (Qwen) and $8/18$ (DeepSeek) test
passes; the edited prompt lifts both models to $17/18$; the manually
selected final implementation passes $18/18$, reaches $100\,\%$ branch
and $100\,\%$ line coverage, and records an $84.95\,\%$ PITest mutation
score. The most important Phase~2 lesson is that, at the integration
level, prompt clarity does more work than model selection. The
combination of branch coverage, hand-crafted mutation guards and
PITest mutation analysis gives three independent angles on test quality.

% =====================================================================
\appendices
\section{Selected HumanEval-X Java Prompts}\label{app:prompts}
% =====================================================================

\begin{table}[!ht]
\renewcommand{\arraystretch}{1.05}
\caption{Phase 1 selected prompts (full list).}
\label{tab:prompts-full}
\centering
\footnotesize
\begin{tabular}{@{}llll@{}}
\toprule
\textbf{Task} & \textbf{Method} & \textbf{Diff.} & \textbf{P2 dep?} \\
\midrule
Java/2  & truncateNumber       & easy  & no \\
Java/3  & belowZero            & easy  & no \\
Java/9  & rollingMax           & easy  & no \\
Java/11 & stringXor            & easy  & no \\
Java/13 & greatestCommonDivisor& easy  & no \\
Java/14 & allPrefixes          & easy  & no \\
Java/22 & filterIntegers       & easy  & no \\
Java/23 & strlen               & easy  & \textbf{yes} \\
Java/24 & largestDivisor       & easy  & no \\
Java/27 & flipCase             & easy  & \textbf{yes} \\
\midrule
Java/0  & hasCloseElements     & moderate & no \\
Java/4  & meanAbsoluteDeviation& moderate & no \\
Java/5  & intersperse          & moderate & no \\
Java/8  & sumProduct           & moderate & no \\
Java/15 & stringSequence       & moderate & no \\
Java/16 & countDistinctChars   & moderate & no \\
Java/18 & howManyTimes         & moderate & \textbf{yes} \\
Java/19 & sortNumbers          & moderate & no \\
Java/25 & factorize            & moderate & no \\
Java/29 & filterByPrefix       & moderate & no \\
\midrule
Java/1  & separateParenGroups  & hard & no \\
Java/6  & parseNestedParens    & hard & no \\
Java/7  & filterBySubstring    & hard & no \\
Java/10 & isPalindrome         & hard & no \\
Java/12 & longest              & hard & no \\
Java/17 & parseMusic           & hard & no \\
Java/20 & findClosestElements  & hard & no \\
Java/21 & rescaleToUnit        & hard & no \\
Java/31 & isPrime              & hard & no \\
Java/39 & primeFib             & hard & no \\
\bottomrule
\end{tabular}
\end{table}

\section{Reproduction Recipe}\label{app:repro}

\begin{lstlisting}[language=bash, caption=Reproducing the Phase 2 BookScan build and PITest report.]
# from Phase2/
mvn -Dmaven.repo.local="$PWD/.m2" clean verify
mvn -Dmaven.repo.local="$PWD/.m2" \
    -P mutation -DskipTests=false test
\end{lstlisting}

The first command runs the 18-test Surefire suite and writes the JaCoCo
report to \texttt{target/site/jacoco/}. The second command runs PITest
and writes the mutation report to \texttt{target/pit-reports/}. The same
flow is automated in \texttt{.github/workflows/phase2-ci.yml}.

% =====================================================================
\section*{Acknowledgment}
% =====================================================================

The two-person group split the project workload along two complementary
tracks; both authors contributed equally and continuously cross-reviewed
each other's deliverables. The same split applied to both phases.
\textbf{Mustafa Eren Ko\c{c}} (150190805) led the \emph{engineering and
automation} track: repository scaffolding, GitHub Actions CI workflows,
the PowerShell + Python~3.11 driver layer, the local \texttt{llama.cpp}
execution pipeline, dataset-to-JUnit adapters, JaCoCo automation, the
Maven and PITest configuration, the \texttt{BookScan} reference
implementation and helper hardening, the prompt-comparison harness and
the Git commit discipline. \textbf{Onat Bar{\i}\c{s} Ercan} (150210075)
led the \emph{quality and analysis} track: 30-prompt selection rationale
and difficulty categorisation, per-task equivalence-partitioning and
boundary-value documents, test-smell audit, hand-crafted mutation-operator
catalogue and \texttt{improvedMutation\ldots} JUnit method authoring, the
\texttt{BookScan} black-box assessment, the original-versus-edited
prompt-strategy design, failure-mode classification, the five-paper
literature review, and the IEEE journal manuscript.

% =====================================================================
\begin{thebibliography}{9}
% =====================================================================

\bibitem{chatunitest}
Z.~Xie, Y.~Chen, C.~Zhi, S.~Deng, and J.~Yin, ``ChatUniTest: a
ChatGPT-based automated unit test generation tool,'' \emph{CoRR},
abs/2305.04764, 2023, doi: 10.48550/arXiv.2305.04764.

\bibitem{metatestgen}
N.~Alshahwan \emph{et al.}, ``Automated unit test improvement using large
language models at Meta,'' in \emph{Companion Proc.\ 32nd ACM Int.\ Conf.
Found.\ Softw.\ Eng. (FSE Companion)}, 2024, pp.~185--196, doi:
10.1145/3663529.3663839.

\bibitem{yang2024}
L.~Yang \emph{et al.}, ``On the evaluation of large language models in
unit test generation,'' in \emph{Proc.\ 39th IEEE/ACM Int.\ Conf.\
Automated Softw.\ Eng. (ASE)}, 2024, pp.~987--999, doi:
10.1145/3691620.3695529.

\bibitem{panta2024}
S.~Gu, N.~Nashid, and A.~Mesbah, ``LLM test generation via iterative
hybrid program analysis,'' \emph{CoRR}, abs/2503.13580, 2025. Accepted
to ICSE 2026, doi: 10.48550/arXiv.2503.13580.

\bibitem{lit_review}
M.~Tasarsu, A.~V.\ Tokmak, and C.~Catal, ``Test case generation using
large language models: a systematic literature review,'' \emph{Cluster
Computing}, vol.~29, no.~4, Springer, 2026, doi:
10.1007/s10586-026-06021-z.

\bibitem{humaneval}
M.~Chen \emph{et al.}, ``Evaluating large language models trained on
code,'' \emph{arXiv preprint arXiv:2107.03374}, 2021.

\bibitem{phase1report}
M.~E.~Ko\c{c} and O.~B.~Ercan, ``Evaluating local code-LLMs on HumanEval-X
Java prompts: semi-agentic generation, improved tests, and mutation-based
guardrails,'' BLG 475E term project, Phase~1 report, Istanbul Technical
University, Apr.\ 2026.

\bibitem{jnose}
J.~Aniche, B.~Alves, and J.~Fonseca, ``JNose: a JUnit test smell detection
tool,'' \emph{Online}: \url{https://jnosetest.github.io/}, accessed
2026-04-25.

\end{thebibliography}

\end{document}
